% Stand 2.08.2026 

%% -----------------------------------
%% |    referencing   |
%% -----------------------------------
% section reference, um z.B. einen Abschnitt der Auswertung in der Diskussion zu referenzieren
\newcommand{\secref}[1]{Abschnitt \labelcref{#1}.(\nameref{#1})}
% um auf einen Stelle zu springen: #1 erwartet das Ziel des Hyperlinks, das kann vieles sein, alles was hyperref erlaubt. e.g. page.number also page.1 springt zu Seite 1. 
% die aktuelle Farbe wird mit colorlet in der variable temp gespeichert und darunter wird eine blaue Linie gezeichnet und danach die temp Farbe wieder eingestellt
\newcommand{\jump}[2]{\hyperlink{#1}{\colorlet{temp}{.}\color{blueurllink}{\underline{\color{temp}#2}}\color{temp}}}


%% -----------------------------------
%% |    textstyles   |
%% -----------------------------------
% Roman Numbers
\newcommand{\RNum}[1]{\uppercase\expandafter{\romannumeral #1\relax}}


%% -----------------------------------
%% |    frequently used math   |
%% -----------------------------------
% the differential operator d not cursive, e.g. in \frac{\d x}{\d t} or in \int dx
\newcommand*\diff{\mathop{}\!\mathrm{d}}      

% partial derivative
\newcommand{\pdb}[2]{\frac{\partial #1}{\partial #2}}

% for A:=B but nice
\newcommand{\defeq}{\vcentcolon=}

% e.g. epsilon up to 0 
\newcommand{\toup}{%
  \mathrel{\nonscript\mkern-1.2mu\mkern1.2mu{\uparrow}}%
}
\newcommand{\todown}{%
  \mathrel{\nonscript\mkern-1.2mu\mkern1.2mu{\downarrow}}%
}

% custom command to make equal sign that can be super- and subscribed
\newcommand{\mytexteq}[2]{%
  \mathrel{%
    \overset{\mbox{\normalfont\tiny\sffamily #1}}{%
      \underset{\mbox{\normalfont\tiny\sffamily #2}}{=}%
    }%
  }%
}
% custom command to make math operations that can be super- and subscribed
\newcommand{\mytextmathop}[3]{%
  \mathrel{%
    \overset{\mbox{\normalfont\tiny\sffamily #1}}{%
      \underset{\mbox{\normalfont\tiny\sffamily #2}}{#3}%
    }%
  }%
}


%% -----------------------------------
%% |    Graphics    |
%% -----------------------------------
% a quick command for including graphics with all necessary vars
% use [angle=phi,width=x] as #2 argument for rotating the image
\newcommand{\xfigure}[5]{\begin{figure}[#1]            
  \centering                                               
  \includegraphics[#2]{#3}                                
  \caption[#4]{#5}
  \label{fig:#3}
\end{figure}}

%Beispiel
%\xfigure{htbp}{width=\textwidth}{filename.png}{Bildbeschreibung fürs \listoffigures}{Bildbeschreibung für unterm Bild \autocite{Skript}}

% for mass including rotated things use this, I put a line/rule over it to distinguish text description from the next image
\newcommand{\xfigurerot}[5]{\begin{figure}[#1]        
\rule{\textwidth}{0.4pt}
\centering
\includegraphics[angle=270,#2]{#3}
\caption[#4]{#5}
\label{fig:#3}
\end{figure}}

% in PAP1 leider manchmal benutzt, und bin zu faul, das richtig zu machen, unterscheidet sich überhaupt nicht zu xfigure, aaußer dass dadurch gekennzeichnet ist, wo pdfs eingefügt werden
% pdfpages fügt leider ganze Seiten ein, d.h. eine Bildunterschrift für eine gescannte PDF ist z.B. nicht möglich
\newcommand{\xfigpdf}[5]{\begin{figure}[#1]
\centering
\includegraphics[#2]{./images/#3}
\caption[#5]{#5}
\label{fig:#3}
\end{figure}}


% I think you'll never need these for PAP, but thanks to https://www.joni.science, you still got them
%% ------------------------------------
%% |    Quantum Mechanics and Math    |
%% ------------------------------------

% this gives one alignable math environment with only one label, e.g for a long error formula
\newenvironment{nalign}{
    \begin{equation}
    \begin{aligned}
}{
    \end{aligned}
    \end{equation}
    \ignorespacesafterend
}

\newcommand{\op}[1]{\operatorname{#1}}                 % to write operators that
                                                       % are not predefined;
                                                       % it's just an abbrev.
                                                       % for the long command

\newcommand{\arr}[2]{\begin{array}{#1}#2\end{array}}   % to create arrays. very 
                                                       % useful in math env

\newcommand{\ket}[1]{\left|#1\right\rangle}           % \ket{X}  ->  |X>
\newcommand{\bra}[1]{\left\langle#1\right|}           % \bra{X}  ->  <X|
\newcommand{\braket}[2]                               % \braket{X}{Y}  ->  <X|Y>
{\left\langle#1 \middle| #2\right\rangle}
\newcommand{\bratenket}[3]                            % \bratenket{X}{Y}{Z}  ->
{\left\langle#1 \middle|\middle| #2 \middle|\middle|  % <X|Y|Z>
#3\right\rangle}
\newcommand{\anglemean}[1]                            % \anglemean{X}  ->  <X>
{\left\langle #1 \right\rangle}                       % \norm{X}  ->  || X ||
\newcommand{\norm}[1]{\left\lVert#1\right\rVert}

\newcommand{\updownarrows}                            % \ket\updownarrows  ->
{\text{\rotatebox[origin=c]{90}{\(\rightleftarrows\)}}} % |↑↓> (cmt is utf8!)
\newcommand{\downuparrows}                            % \ket\updownarrows  ->
{\text{\rotatebox[origin=c]{270}{\(\rightleftarrows\)}}}% |↓↑>
\newcommand{\neswarrows}                              % \ket\neswarrows  ->
{\text{\rotatebox[origin=c]{45}{\(\rightleftarrows\)}}} % |↗↙>
\newcommand{\swnearrows}                              % \ket\swnearrows  ->
{\text{\rotatebox[origin=c]{225}{\(\rightleftarrows\)}}}% |↙↗>

\newcommand{\cre}{c^\dagger}                          % annihalation operator
\newcommand{\anh}{c^{\vphantom{\dagger}}}             % creation operator
\newcommand{\numb}{n^{\vphantom{\dagger}}}            % number operator

\newcommand{\fullstop}{\text{\,.}}                    % fullstop or comma in
\newcommand{\comma}{\text{\,,}}                       % math mode for use
                                                      % after equations
